%======================================================================
\chapter{Discussion}
%======================================================================

\section{Limitations and Challenges}

The development of Mime encountered several significant challenges that ultimately constrained the time available for core optimization and feature expansion. The most substantial hurdle involved 3D asset generation and ARKit rigging. Initially, sourcing a reliable and affordable iOS application capable of generating a high-quality 3D head mesh using the TrueDepth camera proved difficult. While we successfully captured a mesh using the Heges app, the subsequent rigging process introduced severe bottlenecks. Without prior experience or comprehensive guides, we attempted to manually map the ARKit blendshapes to the scanned mesh using Blender and third-party add-ons. Despite a significant time investment, validation tests in tools like glTF Viewer revealed the rigging was structurally flawed and unusable for animation.

This setback forced a necessary pivot to using pre-rigged 3D assets. While these off-the-shelf models functioned reliably and integrated well with our pipeline, they are typically stylized avatars rather than photorealistic human scans. This compromise detracted slightly from our initial objective of creating a highly realistic digital twin.

Beyond the 3D assets, system integration and environment setup presented their own roadblocks. Managing the complex Zoom virtual camera bridge, troubleshooting real-time 3D rendering libraries (specifically pyrender's conversion from .glb to trimesh not saving ARKit blendshape weights), and maintaining parity across both macOS and Windows development environments made it difficult to establish a consistent testing pipeline. Ultimately, troubleshooting these infrastructure and asset pipeline issues consumed a large portion of our development cycle. As a result, less time was available for reducing inference latency and optimizing the core machine learning models, directly reflected in our final results and outlined in the future work section.

\section{Directions for Future Work}

There are several directions to improve Mime into a production-ready system. First, end-to-end latency should be reduced further through tighter streaming orchestration, smaller translation/synthesis models, and improved buffering policies between STT, MT, TTS, and avatar inference. Lower and more stable latency would make turn-taking feel more natural in real conversations.

Second, lip-sync accuracy could be substantially improved by replacing the current audio-driven blendshape model with one that operates directly on the TTS output, to elimiate the temporal misalignment between speech and avatar motion. Grounding predictions on phoneme-level timing as an additional aspect would produce animations that track speech sounds directly rather than approximating them from raw audio.

Third, the blendshape model itself has significant headroom for improvement. The current architecture learns a direct audio-to-coefficient mapping without any explicit notion of viseme state or co-articulation dynamics. Future work could explore sequence-to-sequence models with autoregressive decoding for smoother outputs, or even pretrained audio encoders such as HuBERT or wav2vec 2.0 in place of the custom CNN frontend.

Lastly, the system could integrate the whole setup of the 3D facial scan and rigging process into the pipeline, allowing users to easily create their own personalized avatars without needing to find a pre-rigged model or go through the rigging process themselves. This would make the system more accessible and customizable for users and would more importantly allow for anyone to use a realistic avatar of themselves, which is the ideal use case for this system.

\section{Implications and Impact}

Mime's primary impact is in improving cross-lingual communication quality by addressing both audio and visual coherence, not translation alone. By coupling translated speech with synchronized facial animation, the system targets a common failure mode in multilingual calls: hearing one language while seeing mouth motion from another. Reducing this mismatch can improve perceived naturalness and lower cognitive load for participants.

From a practical perspective, Mime demonstrates that real-time multimodal translation interfaces can run on consumer hardware using modular, API-assisted components. This lowers the barrier for educational teams, small organizations, and independent developers to prototype advanced communication tools without large infrastructure budgets.

More broadly, the project provides a reusable systems template for future telepresence applications where timing alignment between modalities is critical, including accessibility support, multilingual tutoring, and remote collaboration.

\section{Conclusions}

Mime demonstrates that real-time speech translation with synchronized facial animation is achievable on commodity hardware. The pipeline successfully integrates STT, MT, TTS, and audio-driven blendshape prediction into a single working system, validating the core idea that translation and lip-sync are better solved together than separately.

There is clear room for improvement in latency, model accuracy, and robustness, but the current implementation serves as a solid foundation for future work in cross-language video communication.